\chapter{The vision pipeline, measured on the real cameras}
\label{ch:vision}

\emph{Measured on four physical cameras, 2026-08-30. This chapter and
\cref{sec:hardware-session} are the only places in this thesis where the
numbers come from hardware in the laboratory rather than from the simulated
platform.}

\bigskip

``The vision works'' is not a claim a thesis can print, and until this session
it was the only claim available: the perception stages of \cref{ch:perception}
had been exercised against rendered scenes, where the answer is known because
the scene was constructed, and against no real camera at any point. This
chapter reports what happens when the same code is pointed at four physical
cameras and every stage is required to show its working.

\section{Method: one figure and one formula per stage}
\label{sec:vision-method}

The pipeline was decomposed into thirty-three stages, each of which is one
operation with one formula. Every stage draws itself on the frame it was
given, states the operation it performed, and reports the quantities it
produced \emph{with units}; the numbers are written to a machine-readable
table at the same time as they are burned into the figure, so a figure and the
table beneath it cannot disagree. The run of 2026-08-30 produced
\num{723} such measurements.

Three properties of the instrument matter more than any single number it
produced.

\textbf{A stage that cannot run says so, and says what would fix it.} There
are three outcomes, not two. A stage \emph{succeeds}; a stage
\emph{refuses}, and the refusal names its cause and its remedy; or a stage is
\emph{absent by design}, which is drawn differently again. The webcam has no
depth sensor, so its fifteen depth stages are by-design absences and are not
faults. Rendering a by-design absence identically to a defect is a failure
mode this project has recorded before, and it costs a reader the ability to
tell a limitation from a bug.

\textbf{Nothing in the file can move the arm.} The program contains no
publisher, no service client, no Kortex session and no trajectory. Its only
possible traffic is a subscription. This is a structural property rather than
a policy, and it is what allows the pipeline to be exercised on live hardware
without a safety case.

\textbf{The arithmetic is checked against a constructed frame whose answer is
known.} A self-test runs every geometric stage against a synthetic depth image
whose plane, object heights and object widths were chosen rather than
measured, so that a stage returning the wrong answer fails before any real
camera is involved. This is the standing rule of \cref{ch:instrument} applied
to a new instrument at the point of building it rather than after it produced
a surprising number.

\subsection{The geometry, written out}

Depth arrives as a range image and is lifted to a point cloud by the pinhole
model, using intrinsics read from the device rather than assumed:

\begin{equation}
  \label{eq:deproject}
  \mathbf{P}(u,v) =
  \left(
    \frac{u - c_x}{f_x}\,Z(u,v),\;
    \frac{v - c_y}{f_y}\,Z(u,v),\;
    Z(u,v)
  \right)^{\!\top} .
\end{equation}

The support plane is fitted by random sample consensus. Three points are drawn,
the candidate plane is
$\mathbf{n} = (\mathbf{b}-\mathbf{a}) \times (\mathbf{c}-\mathbf{a})$
normalised, with offset $d = -\mathbf{n}\cdot\mathbf{a}$, and it is scored by
the number of points within a tolerance $\tau = \SI{6}{\milli\metre}$:

\begin{equation}
  \label{eq:ransac}
  S(\mathbf{n}, d) =
  \bigl|\bigl\{\,\mathbf{P} : |\mathbf{n}\cdot\mathbf{P} + d| < \tau \,\bigr\}\bigr| ,
  \qquad
  \text{best of \num{400} draws.}
\end{equation}

Height above the fitted surface is then the signed residual, and it is what
separates an object from the surface it stands on:

\begin{equation}
  \label{eq:height}
  h(\mathbf{P}) = \mathbf{n}\cdot\mathbf{P} + d .
\end{equation}

Two quantities are reported for every plane fit and both are needed. The
residual root mean square says how flat the fitted set is; the inlier
\emph{fraction} says whether the fitted set is the scene. A plane fits three
points perfectly, so a small residual on its own establishes nothing, and
\cref{sec:vision-scene} is a case where the two diverge sharply.

\section{The orientation defect}
\label{sec:orientation}

The first thing the instrumentation found was in the instrumentation's own
input, which is the outcome the standing rule predicts.

Three separate places in the software decided independently whether the scene
depth camera is mounted upside down. Two rotated its frames by
\SI{180}{\degree} and carried the comment that the camera is inverted; the
third, the live node that publishes to the rest of the system, rotated
nothing. Nobody had measured it, so the recorded figures and the live topics
disagreed and neither side could tell.

It was settled by measurement rather than by inspection, using a reference
that needs no convention: the HD webcam watches the same scene from the same
side of the room and is rotated by nothing. Each frame is reduced to a
$64\times 64$ grey signature, standardised so that a brighter camera does not
score differently from a darker one, and correlated against the reference in
both orientations. Over four independent captures the un-rotated RealSense
frame agrees with the webcam at $r = \numrange{0.66}{0.70}$ and the rotated
one at $r = \numrange{0.27}{0.34}$: the same verdict every time, a factor of
\num{2.4} apart. \textbf{The camera is mounted the right way up}, and every
scene figure recorded before the correction was inverted.

The correction is not a rotation of the image. A \SI{180}{\degree} rotation
maps pixel $(u,v)$ to $(W\!-\!1\!-\!u,\; H\!-\!1\!-\!v)$, and the principal
point is a pixel, so it must move with it:

\begin{equation}
  \label{eq:rot180}
  c_x' = W - 1 - c_x, \qquad c_y' = H - 1 - c_y .
\end{equation}

Rotating the image and leaving the intrinsics alone places every deprojected
ray on the wrong side of the optical axis. On this device the principal point
sits \SI{8.85}{px} from the image centre, so the error is \SI{17.7}{px} of
apparent principal-point displacement, or \SI{29}{\milli\metre} at
\SI{1}{\metre} --- small, real, and completely invisible in a picture. The
orientation is now stored as a measured configuration with the evidence that
set it, one function applies it to the image and the intrinsics together so
that half a rotation cannot be applied, and all three consumers read the same
file. Undoing the rotation restored the intrinsics to the device's own
reported values of $c_x = \num{318.87}$, $c_y = \num{231.22}$ exactly, which
is an independent confirmation that \cref{eq:rot180} was applied in the right
direction.

\subsection{What the defect did and did not corrupt}

Because the rotation was applied \emph{consistently} to the image and the
intrinsics, substituting \cref{eq:rot180} into \cref{eq:deproject} negates the
$X$ and $Y$ components and leaves $Z$ untouched. The recovered point cloud was
therefore an exact reflection of the true one through the optical axis, and
every distance, extent, angle and residual is preserved under a reflection. No
metric quantity should have changed at all, and this was checked rather than
assumed: the deprojected extents are exactly negated
($X \in [\num{-3.477}, \num{1.141}]$ became
$X \in [\num{-1.141}, \num{3.477}]$) and the centroid components change sign
and nothing else.

Two classes of stage changed anyway, and the reason each changed is worth
recording.

\textbf{The learned segmenter is not rotation-invariant.} Given the same scene
the right way up it returned \num{96} regions rather than \num{92}, with a
different largest region and a different median area. A convolutional
segmenter has no reason to be invariant to a \SI{180}{\degree} rotation and
was not trained to be, so this is expected behaviour and not a defect --- but
it means that a geometric pipeline can be immune to an orientation error while
the learned component sitting beside it is not, and only the learned component
will show it.

\textbf{The plane fit changed because its support is weak, not because the
geometry changed.} \Cref{eq:ransac} draws candidate triples in raster order,
and rotating the image reverses that order, so a different sequence of
candidates is examined. On the wrist cameras, where the best plane holds
\SI{74.0}{\percent} of the points, this makes no difference. On the scene
camera, where the best plane holds \SI{3.8}{\percent}, it selects a different
plane entirely: \SI{2.68}{\metre} from the camera before the correction and
\SI{4.35}{\metre} after. That is not a consequence of the orientation bug. It
is a consequence of pointing a plane-finder at a room, and the orientation
change merely made it visible.

\section{What the cameras measured}
\label{sec:vision-results}
\label{sec:vision-scene}

\begin{table}[htbp]
  \centering
  \small
  \caption[Stage outcomes by camera]{Outcomes of \num{33} stages on each of
  four cameras. \emph{By design} is the webcam's depth stages, which are
  absent because the sensor is absent; they are counted separately from
  refusals so that a limitation is never reported as a fault. \textbf{No
  stage crashed on any camera.}}
  \label{tab:vision-stages}
  \begin{tabular}{@{}lcccl@{}}
    \toprule
    Camera & Result & Refused & By design & What it is \\
    \midrule
    left wrist  & 28 & 5 & 0
      & Kinova vision module, RGB-D, on the arm \\
    right wrist & 29 & 4 & 0
      & Kinova vision module, RGB-D, on the arm \\
    scene depth & 27 & 6 & 0
      & RealSense D435i across the room \\
    scene colour & 12 & 8 & 13
      & HD webcam, colour only, never calibrated \\
    \midrule
    \textbf{total} & \textbf{96} & \textbf{23} & \textbf{13} & of \num{132} \\
    \bottomrule
  \end{tabular}
\end{table}

\begin{table}[htbp]
  \centering
  \small
  \caption[Geometry recovered from each depth camera]{Geometry recovered from
  each depth camera in one frame. The inlier fraction is the column that
  decides whether the row below it means anything: the wrist cameras identify
  a work surface and the scene camera does not, which is the honest
  consequence of pointing a plane-finder at a room rather than at a table.}
  \label{tab:vision-geometry}
  \begin{tabular}{@{}lccc@{}}
    \toprule
    & left wrist & right wrist & scene depth \\
    \midrule
    pixels returning a depth      & \SI{89.4}{\percent} & \SI{80.5}{\percent} & \SI{14.5}{\percent} \\
    median range                  & \SI{2.222}{\metre} & \SI{1.838}{\metre} & \SI{4.706}{\metre} \\
    points inside the working band & \num{54789} & \num{42779} & \num{38526} \\
    colour pixels given a depth   & \SI{64.4}{\percent} & \SI{58.8}{\percent} & aligned in hardware \\
    \midrule
    plane inliers                 & \SI{74.0}{\percent} & \SI{52.8}{\percent} & \SI{3.8}{\percent} \\
    plane residual, RMS           & \SI{1.43}{\milli\metre} & \SI{1.59}{\milli\metre} & \SI{3.41}{\milli\metre} \\
    surface extent                & $\num{0.579}\times\num{0.423}$ & $\num{0.428}\times\num{0.467}$ & $\num{1.012}\times\num{1.762}$ \\
    surface distance              & \SI{0.278}{\metre} & \SI{0.290}{\metre} & \SI{4.352}{\metre} \\
    \midrule
    objects posed in three dimensions & 5 & 12 & 12 \\
    nearest object                & \SI{0.539}{\metre} & \SI{0.387}{\metre} & \SI{2.609}{\metre} \\
    graspable within the jaw      & 0 & 0 & 0 \\
    \bottomrule
  \end{tabular}
\end{table}

\begin{figure}[htbp]
  \centering
  \begin{subfigure}[b]{0.49\linewidth}
    \centering
    \includegraphics[width=\linewidth]{vision/scene_rs_colour.png}
    \caption{Colour, the right way up.}
    \label{fig:rs-colour}
  \end{subfigure}\hfill
  \begin{subfigure}[b]{0.49\linewidth}
    \centering
    \includegraphics[width=\linewidth]{vision/scene_rs_depth.png}
    \caption{The aligned depth of the same frame. Magenta is no return.}
    \label{fig:rs-depth}
  \end{subfigure}
  \caption[The scene camera after the orientation correction]{The scene
  camera after the correction of \cref{sec:orientation}. The depth panel is
  the reason the row for this camera in \cref{tab:vision-geometry} reads as it
  does: at a median range of \SI{4.706}{\metre} only \SI{14.5}{\percent} of
  pixels return anything, and what does return is the far wall and the floor
  rather than a work surface. A stereo depth camera across a room is not a
  small version of a wrist camera.}
  \label{fig:rs-pair}
\end{figure}

\Cref{tab:vision-geometry} contains one result and one non-result, and it is
worth being explicit about which is which.

\textbf{The wrist cameras measure a work surface to about
\SI{1.5}{\milli\metre}.} A plane holding \SI{74.0}{\percent} of the returned
points at \SI{1.43}{\milli\metre} root mean square, at a working distance of
\SI{0.278}{\metre}, is the quantity the autonomous pick path in
\cref{ch:perception} depends on, and it is close to the \SI{0.54}{\milli\metre}
this project recorded for the same sensor class against a rendered scene. This
is the first evidence in the project that the surface-fitting stage behaves on
a real table as it does on a constructed one.

\textbf{The scene camera does not identify a support plane, and the figure
says so rather than reporting the number quietly.} Its best plane holds
\SI{3.8}{\percent} of the points. Every quantity that depends on it --- height
above the surface, object heights, three-dimensional boxes --- therefore rests
on a plane that one point in twenty-six agrees with, and none of those numbers
should be quoted. The instrument reports them because suppressing them would
hide the inlier fraction that condemns them; the interpretation is stated here
instead.

\subsection{The two defects the first runs exposed}

The first attempts at this session produced results that were wrong in ways
that only a stage-by-stage instrument would have separated from a broken
camera, and both were in the software.

\emph{A gate belongs to the thing being measured.} The depth working band
$\SIrange{0.08}{1.5}{\metre}$ is the pick path's own, and is correct for a
camera on the hand. Applied unchanged to a camera across the room, whose
nearest return is \SI{0.741}{\metre} and whose median is \SI{3.84}{\metre}, it
retained \textbf{\num{0} of \num{85061} returns}, and sixteen stages then
refused for want of depth that the camera had in abundance. The band is now
declared per source and every figure states the band it used; the same camera
now yields \num{38526} points.

\emph{Two stages disagreed about the same frame, and the disagreement was the
evidence.} The stage that asks what is standing on the surface reported
nothing, while the stage that poses objects in three dimensions --- on the
same frame --- reported six objects, the nearest at \SI{0.536}{\metre}. Both
cannot be true. The cause was that depth at $\num{480}\times\num{270}$
forward-projected into colour at $\num{1280}\times\num{720}$ lands one sample
in roughly every seven colour pixels, so only \SI{7.22}{\percent} of colour
pixels carried a depth. The above-plane mask was a dot pattern: every
connected component fell below the area floor, so the stage that needs
connectivity saw nothing while the stage that works per mask was unaffected.
Splatting each projected sample over the gap it spans, with the radius derived
from the resolution ratio rather than chosen, raised coverage to
\SI{64.4}{\percent} and \SI{58.8}{\percent}
(\cref{tab:vision-geometry}). A contradiction between two stages was worth
more than either stage's output.

\subsection{What did not work, and why that is an answer}

\textbf{No object in any frame could be grasped.} The parallel-jaw planner
refused on every camera, and the refusal carries its arithmetic: the narrowest
axis of the nearest object measures \SI{144}{\milli\metre} against a usable
jaw opening of \SI{75}{\milli\metre}. That is a statement about the scene, not
a failure of the planner, and it is the same geometric constraint that
\cref{sec:cells} imposes on the study's carry tasks.

\textbf{The open-vocabulary detector returned nothing on any camera.} It was
given the single class \emph{green cube} and found no instance of it at a
confidence floor of \num{0.05}. It was not asked for people or furniture, so
this is the measured behaviour of that model on this rig with that prompt
rather than a misconfiguration; the classical colour and geometry stages found
the same objects the model did not.

\textbf{No fiducial marker decoded in any frame}, because none was in the
scene. A marker either decodes or it does not, which is the property that
makes it a poor thing to depend on and a good thing to check.

\textbf{The colour scene camera has never been calibrated}, so eight of its
stages refuse rather than assume that the principal point is the image centre.
On the depth camera the principal point is \SI{8.85}{px} from centre, which
would be \SI{14.7}{\milli\metre} of error at \SI{1}{\metre} if assumed away.
The refusal is the correct behaviour and the calibration is a half-hour task
that has not been done.

\textbf{The robot is not detected in any image, and no stage claims that it
is.} Nothing in this system finds an arm in a picture. The arm's pose comes
from its encoders through the kinematic description, and drawing it into a
camera image requires the hand-eye transform that \cref{ch:motion} identifies
as the largest unmeasured term in the positioning budget.

\subsection{A caption that could not be written}

The arms were to be parked at the pick pose for this session so that the
figures could be captioned as such. The instrument measures the pose rather
than assuming it, and found the left arm \SI{10.67}{\degree} and the right
\SI{13.53}{\degree} from the reference configuration, against a tolerance of
\SI{2}{\degree}. Earlier runs in the same session were \SIrange{66}{82}{\degree}
out.

None of these figures may therefore be captioned ``at the pick pose'', and
none is. This is the check working rather than failing: a figure captioned
with a pose nobody measured is worth nothing, and nothing else in the session
would have noticed. It is the same failure mode as \cref{sec:orientation} in a
different variable --- a property everybody believed, that nobody had put a
number on.

\section{What this chapter changes}

Three things are true after this session that were not true before it.

The perception geometry has been exercised on real sensors and the surface fit
holds to \SI{1.43}{\milli\metre} at working distance, which is the term the
pick path depends on and which had previously only been characterised against
a renderer.

The scene camera is a different instrument from the wrist cameras, not a
distant one, and the numbers say by how much: an eighth of the depth returns,
a twentieth of the plane support, and a scene with no dominant plane in it.
Any autonomous behaviour that expects a support surface must take it from a
wrist camera.

And a constant that three separate files had asserted and none had measured
was wrong. It was found in the first hour of pointing an instrumented pipeline
at the hardware, it was invisible to every metric stage by construction, and
the only thing that would ever have shown it is a second camera looking at the
same room. It is the twenty-sixth entry in \cref{ch:instrument}, and the first
whose mechanism is a constant asserted in three places and measured in none.
